\section{Intrinsic common-port atlases for surface discovery}
\label{sec:intrinsic-port-atlases}

Report~113 supplies an objective-independent projective atlas for a
collection of normal-coordinate blocks sharing one common port.  Its
exponent depends on the number of port coordinates.  The new maintained
preprocessor \path{fastunknot/port_subspace.py} replaces that ambient
dimension by the dimension of the intersection of the blocks' linear
port images.  This is a polynomial exact reduction of the discovery
query.  The atlas remains a research backend, accessed through
\path{projective_port_research/common.py}; the ordinary recognizer's
dispatch is unchanged.

\subsection{The common-port model and the inherited atlas theorem}

Block $i$ has coordinates $x_i\geq0$, rational matching equations
$A_i x_i=0$, and pairwise coordinate incompatibilities.  Its port matrix
$P_i$ is entrywise nonnegative with $d$ rows.  Feasible tuples satisfy
$P_i x_i=y$ for every $i$.  The query asks whether
\[
 c_0^T y+\sum_i c_i^T x_i>0,
\]
possibly with specified coordinates set to zero.  All blocks share this
\emph{same} $y$; this is not a theorem for arbitrary chains or trees of
different seams.  Cross-block incompatibilities must be accounted for by
the source compiler, not silently discarded.

For each block enumerate its admissible extreme rays.  A positive
zero-port ray answers the query immediately; other zero-port summands
can be removed without decreasing the objective.  Normalize every
remaining ray by the sum of its port coordinates.  At fixed normalized
$y$, an optimal compatible local mixture uses at most $d$ linearly
independent port columns.  These mixtures define simplices in the
standard $(d-1)$-simplex, with affine objective on each chosen local
piece.  The common refinement has an objective-independent vertex atlas.
For $d=2$ the union of local normalized ray images suffices.  For larger
$d$, local ray images alone do not suffice: two medians can meet at a
new interior point.  Report~113 constructs the required refinement by
adjoining the coordinate frame to each local point set, spanning
hyperplanes with $(d-1)$-subsets, and intersecting $(d-1)$ such
hyperplanes with the normalization hyperplane.  Extending independent
simplex columns by coordinate vectors shows that every simplex face is
in the arrangement.  An affine objective attains its maximum at a
vertex of the relevant common cell.  Keeping all local mixtures at
each vertex also preserves later coordinate-zero queries.

With $R$ enumerated local rays, the inherited bound after enumeration is
\[
 \operatorname{poly}(\Lambda,d)(R+d+1)^{d^2-d+1},
\]
where $\Lambda$ includes the binary source and query encoding.  Complete
support enumeration costs $2^{n_i}\operatorname{poly}(\Lambda)$ per block;
the alternative zero-set enumeration examines
$\binom{n_i}{h_i-1}$ candidates when $h_i=\dim\ker A_i>0$.
These are actual enumeration costs, not unit-cost oracles.

\subsection{A polynomial intrinsic-dimension reduction}

\begin{theorem}[Common-image restriction]
\label{thm:intrinsic-common-port}
Let $W_i=P_i\ker A_i\subseteq\mathbb Q^d$ and
$W=\bigcap_i W_i$, of dimension $r$.  In polynomial bit time one can
construct new block equations and nonnegative port matrices with $r$
rows, when $r>0$, such that synchronized feasible coordinate tuples are
exactly the original ones.  Coordinates, incompatibilities, and all
coordinate-zero restrictions are preserved.  Every linear objective
transforms exactly.  For $r=0$, one dummy identically zero port row
suffices.
\end{theorem}

\begin{proof}
Compute bases of $\ker A_i$, their images under $P_i$, and their
annihilators.  The sum of these annihilators is $W^\perp$; let its row
basis be $H$.  Replace each local equation matrix by a row basis of
\[
 A'_i=\begin{pmatrix}A_i\\ HP_i\end{pmatrix}.
\]
Every original synchronized tuple has common port in every $W_i$,
hence satisfies these added equations.  Conversely they merely require
the individual port images to belong to $W$.

Put a row basis $T$ of $W$ in reduced row echelon form and let $J$ be
its $r$ pivot columns.  Then $T_{:,J}=I_r$.  Projection to coordinates
$J$ is injective on $W$, and a port with selected coordinates $z$ has
unique full lift $y=T^Tz$.  Set $P'_i=(P_i)_{J,:}$.  These matrices
remain entrywise nonnegative, even if $T$ contains negative entries.
Equality of the reduced ports is equivalent to equality of the full
ports, since all full ports now lie in $W$.  No inequalities implied by
the unselected coordinates are lost: the original nonnegative variables
and equations remain.  Replace the master objective by $c'_0=Tc_0$;
then $(c'_0)^Tz=c_0^Ty$ exactly.  Leave all local objectives unchanged.
If $r=0$, the added equations force every full port to vanish; a dummy
zero row and zero master objective implement this case.

The construction uses a polynomial number of rational matrix products
and eliminations of polynomial size.  The bit lengths of the resulting
bases are polynomially bounded by minors after denominator clearing;
there is no iterative enumeration hidden in the construction.  This
also bounds the added constraints and lifting matrix in the binary
input length.
\end{proof}

The word ``linear'' is essential: $W_i$ is a linear image of a kernel,
not a claim about the span of a compatible nonnegative cone.  Ignoring
nonnegativity at this preparatory step yields a safe overestimate, which
may leave further reductions undiscovered.  No claim of minimal
possible port dimension is made.

\begin{corollary}[Intrinsic atlas exponent]
After recomputing rays of the restricted blocks, let their total number
be $R'$ and their nullities be $h'_i\leq h_i$.  For $r\geq1$ the atlas
bound becomes
\[
 \operatorname{poly}(\Lambda,r)(R'+r+1)^{r^2-r+1},
\]
in addition to the polynomial preprocessor and the new local
enumerations.  Thus the logarithmic common-port hypothesis in the
conditional quasi-polynomial result can be imposed on $r$, even when
the original $d$ is larger.
\end{corollary}

One must \emph{recompute} local rays.  For example, a block with no
equations and identity two-coordinate port has rays $(1,0)$ and
$(0,1)$.  Synchronizing with a one-variable block of port $(u,u)$ adds
$x_1=x_2$ to the first block.  Its new ray $(1,1)$ is neither old ray.
There is no asserted inequality $R'\leq R$.  Small intrinsic dimension
alone therefore does not bound the entire algorithm.  If both local
coordinate counts and $r$ are $O(\log N)$, with polynomial binary
source length, the displayed costs give $2^{O(\log^3N)}$; bounded local
nullity provides a separate way to control ray enumeration.

The projective normalization also changes: $\sum_j z_j=1$ need not
equal $\sum_j y_j=1$.  The feasible homogeneous tuples and objective
signs agree, but normalized optimal values generally do not.  The
research adapter returns a lifted full port, the original coordinate
vectors, and the objective evaluated on those vectors.  It does not
compare differently normalized optima.

\subsection{Source recovery and the unresolved geometric hypotheses}

An assembled positive vector need not be extreme, even after primitive
integer scaling.  The report's source recovery normalizes its coordinate
sum and descends its support: choose a nonzero vector in the supported
matching kernel with coordinate sum zero, orient it so the objective
does not decrease, and move until a positive coordinate vanishes.
There are at most $n-1$ reductions.  Integer null vectors have
polynomially bounded entries from minors, and each step multiplies a
common denominator by at most one such entry; hence the entire descent
has polynomial bit complexity.  It preserves positivity, conflicts,
and forced zeros.  The final primitive matching ray has one-dimensional
support kernel, which the independent source verifier checks directly.

For authenticated standard normal coordinates, such a primitive
admissible extreme ray is connected: every component vector lies on
the same extreme ray, and integrality plus primitivity makes its scale
a positive integer.  This does not itself prove essential boundary,
diagram provenance, or a useful hierarchy reduction.  Positive
normal-surface candidates still require the existing independent native
disc checks.  Complete negative atlas answers concern only the supplied
algebraic query; a resource limit is inconclusive.

Neither report~113 nor this reduction proves that all required native
discovery problems possess suitable small blocks or small intrinsic
common-port dimension.  In particular, a constant dimension on each
seam of an arbitrary chain is insufficient; the report encodes a
binary-gain product-choice problem in such chains.  Likewise deleting
$s$ vertices from a connected degree-$\Delta$ tetrahedral adjacency
graph leaves at most $\Delta s$ private components.  If their sizes are
at most $b$, then $t\leq s+\Delta sb$, obstructing simultaneous
logarithmic raw separator and raw component sizes for large $t$.
Algebraic compression must be proved, not inferred from a drawing of
small separators.  Useful-surface discovery, marked successor repair,
and aggregate hierarchy costs remain open obligations for the general
quasi-polynomial recognizer.
